\documentclass[11pt,a4paper]{article}
\usepackage[utf8]{inputenc}
\usepackage[T1]{fontenc}
\usepackage[french]{babel}
\usepackage[margin=2.2cm]{geometry}
\usepackage{booktabs,tabularx,array}
\usepackage{amssymb,amsmath}
\usepackage{enumitem}
\usepackage{graphicx}
\usepackage{xcolor}
\usepackage{hyperref}
\hypersetup{colorlinks=true,linkcolor=blue!50!black,urlcolor=blue!50!black}
\setlist{nosep,leftmargin=*}
\setlength{\tabcolsep}{4pt}

\title{\textbf{CareStep} --- Labs : Pr\'eparation, Mod\'elisation, \'Evaluation\\
\large Fyras $\mid$ Projet ML avanc\'e --- Dataset Diabetes 130-US Hospitals}
\author{Fyras}
\date{}

\begin{document}
\maketitle

% ─────────────────────────────────────────────────────────────────────────────
\section{Vue d'ensemble du pipeline}

Le pipeline CareStep transforme les donn\'ees brutes du dataset UCI \#296 en un score de risque calibr\'e, puis en un plan d'actions personnalis\'e pour chaque patient.

\medskip
\noindent
\fbox{\parbox{0.95\textwidth}{\centering
\textbf{Donn\'ees brutes} $\;\to\;$
\textbf{Pr\'eparation} $\;\to\;$
\textbf{Mod\'elisation} $\;\to\;$
\textbf{\'Evaluation} $\;\to\;$
\textbf{Tuning} $\;\to\;$
\textbf{SHAP $\to$ Plan d'action}
}}

\medskip
Chaque \'etape correspond \`a un script Python (un ``Lab'') ex\'ecutable ind\'ependamment, chaque Lab chargeant les r\'esultats du pr\'ec\'edent.

% ─────────────────────────────────────────────────────────────────────────────
\section{Lab 1 --- Pr\'eparation des donn\'ees}

\subsection{Chargement et nettoyage}

\begin{itemize}
  \item Dataset : 101\,766 encounters, 50 colonnes, 71\,518 patients uniques.
  \item \textbf{Exclusion} des sorties d\'ec\`es/palliatifs (\texttt{discharge\_disposition\_id} $\in \{11,13,14,19,20,21\}$) $\to$ 2\,423 lignes retir\'ees.
  \item \textbf{Exclusion} du genre \texttt{Unknown/Invalid} (r\`egle 2 du protocole) $\to$ 3 lignes retir\'ees $\to$ \textbf{99\,340} encounters restants.
\end{itemize}

\subsection{Construction de la cible}

$$y = \mathbb{1}[\texttt{readmitted} = \text{``<30''}]$$

Pr\'evalence : $\approx 11{,}4\,\%$ de positifs (classe d\'es\'equilibr\'ee).

\subsection{Les 12 features patient}

On utilise uniquement des features qu'un patient peut lui-m\^eme r\'epondre \`a la sortie :

\begin{center}
\small
\begin{tabularx}{\textwidth}{@{}clX@{}}
\toprule
\# & \textbf{Feature} & \textbf{Transformation} \\
\midrule
1 & \texttt{age\_ord} & Tranche d'\^age $\to$ ordinal (0--9) \\
2 & \texttt{gender\_bin} & Homme = 1, Femme = 0 \\
3 & \texttt{inpatient\_cap} & Hospitalisations pass\'ees, cap \`a 3 \\
4 & \texttt{emergency\_cap} & Urgences pass\'ees, cap \`a 3 \\
5 & \texttt{outpatient\_bkt} & Consultations pass\'ees, buckets (0/1/2/3) \\
6 & \texttt{time\_in\_hospital} & Dur\'ee du s\'ejour (1--14) \\
7 & \texttt{diag\_group} & Code ICD-9 regroup\'e en 7 familles \\
8 & \texttt{num\_medications} & Nombre de m\'edicaments \\
9 & \texttt{insulin} & Changement insuline (Up/Down/Steady/No) \\
10 & \texttt{change} & Autre m\'edicament chang\'e (Ch/No) \\
11 & \texttt{a1c\_tested} & Test HbA1c effectu\'e (0/1) \\
12 & \texttt{discharge\_group} & Destination (home/home\_with\_help/facility) \\
\bottomrule
\end{tabularx}
\end{center}

Les 7 familles de diagnostic : circulatory (29{,}9\,\%), other (21\,\%), respiratory (14\,\%), infection/renal (10{,}1\,\%), digestive (9{,}4\,\%), diabetes (8{,}7\,\%), injury (6{,}9\,\%).

\textbf{Pi\`ege corrig\'e --- \texttt{a1c\_tested}.} Dans le CSV, la valeur texte \texttt{None} de \texttt{A1Cresult} signifie ``test non fait'', mais pandas la lit comme une valeur manquante (\texttt{NaN}). Le test \texttt{A1Cresult != "None"} donnait donc 1 pour \emph{toutes} les lignes (feature constante, inutile). On utilise \texttt{A1Cresult.notna()} : 16{,}9\,\% des s\'ejours ont un test HbA1c, 83{,}1\,\% non.

\subsection{Split anti-fuite}

\textbf{GroupShuffleSplit} sur \texttt{patient\_nbr} (80/20) : un patient n'appara\^it que dans un seul ensemble. Cela \'evite la fuite d'information o\`u le mod\`ele ``reconna\^it'' un patient d\'ej\`a vu.

\begin{center}
\small
\begin{tabular}{@{}lrl@{}}
\toprule
\textbf{Ensemble} & \textbf{Encounters} & \textbf{Patients} \\
\midrule
Train & 79\,567 & 55\,989 \\
Test & 19\,773 & 13\,998 \\
\bottomrule
\end{tabular}
\end{center}

Patients en commun entre train et test : \textbf{0}.

\subsection{Preprocesseur sklearn}

Un \texttt{ColumnTransformer} transforme les 12 features en 23 colonnes num\'eriques (passthrough pour les num\'eriques, OneHotEncoder pour les cat\'egorielles). Un deuxi\`eme preprocesseur encode \emph{toutes} les features (98 colonnes) pour le mod\`ele de r\'ef\'erence. \texttt{race} n'y figure pas : elle sert uniquement \`a l'audit d'\'equit\'e, jamais en entr\'ee d'un mod\`ele.


% ─────────────────────────────────────────────────────────────────────────────
\section{Lab 2 --- Mod\'elisation}

\subsection{Strat\'egie : 3 \'etapes progressives}

\begin{enumerate}
  \item \textbf{R\'egression logistique} (baseline) --- 12 features patient, \texttt{class\_weight="balanced"}.
  \item \textbf{LightGBM patient} --- 12 features, \texttt{scale\_pos\_weight} = 7{,}8 pour compenser le d\'es\'equilibre.
  \item \textbf{LightGBM FULL} --- toutes les features (98 colonnes) comme r\'ef\'erence.
\end{enumerate}

\subsection{R\'esultats de base (avant tuning)}

\begin{center}
\begin{tabular}{@{}lcc@{}}
\toprule
\textbf{Mod\`ele} & \textbf{PR-AUC} & \textbf{ROC-AUC} \\
\midrule
LogReg (12 features) & 0{,}190 & 0{,}659 \\
LightGBM (12 features) & 0{,}202 & 0{,}653 \\
LightGBM FULL (r\'ef\'erence) & 0{,}208 & 0{,}659 \\
\bottomrule
\end{tabular}
\end{center}

\textbf{Ratio patient/full (mod\`eles bruts)} = $0{,}202 / 0{,}208 = 96{,}8\,\%$ --- on perd tr\`es peu de performance en n'utilisant que 12 features.

\subsection{Calibration}

Le LightGBM produit des scores, pas des probabilit\'es fiables. On applique une calibration isotonique :
\begin{enumerate}
  \item S\'eparer le train en 80\,\% (entra\^inement) et 20\,\% (calibration) avec \texttt{GroupShuffleSplit}.
  \item Entra\^iner un LightGBM frais sur les 80\,\%.
  \item Ajuster une courbe isotonique sur les 20\,\% (\texttt{CalibratedClassifierCV}, \texttt{cv="prefit"}).
\end{enumerate}

R\'esultat : le Brier score passe de 0{,}222 \`a \textbf{0{,}097} --- les probabilit\'es pr\'edites correspondent mieux \`a la r\'ealit\'e.

\subsection{Niveaux de risque}

Deux seuils simples bas\'es sur la pr\'evalence $p \approx 11\,\%$ :

\begin{center}
\begin{tabular}{@{}lcl@{}}
\toprule
\textbf{Niveau} & \textbf{R\`egle} & \textbf{Action} \\
\midrule
Low & $\hat{p} < p$ & Plan standard \\
Medium & $p \le \hat{p} < 2p$ & Plan + suivi rapproch\'e \\
High & $\hat{p} \ge 2p$ & Plan + alertes \'equipe soignante \\
\bottomrule
\end{tabular}
\end{center}


% ─────────────────────────────────────────────────────────────────────────────
\section{Lab 3 --- \'Evaluation}

\subsection{M\'etriques}

\begin{center}
\begin{tabular}{@{}lccc@{}}
\toprule
\textbf{Mod\`ele} & \textbf{PR-AUC} & \textbf{ROC-AUC} & \textbf{Brier} \\
\midrule
LogReg (12f) & 0{,}190 & 0{,}659 & 0{,}230 \\
LightGBM (12f) & 0{,}202 & 0{,}653 & 0{,}222 \\
LightGBM calibr\'e (12f) & 0{,}189 & 0{,}647 & \textbf{0{,}097} \\
LightGBM FULL (r\'ef.) & 0{,}208 & 0{,}659 & 0{,}211 \\
\bottomrule
\end{tabular}
\end{center}

\textbf{Ratio PR-AUC patient calibr\'e/full = 0{,}189 / 0{,}208 = 90{,}9\,\%} $\ge$ cible de 85--90\,\% $\checkmark$ \\
(Le mod\`ele d\'eploy\'e est le mod\`ele calibr\'e ; la calibration isotonique co\^ute un peu de PR-AUC, d'o\`u 90{,}9\,\% contre 96{,}8\,\% pour les mod\`eles bruts.)

\subsection{Recall par niveau de risque}

\begin{center}
\begin{tabular}{@{}lrrr@{}}
\toprule
\textbf{Niveau} & \textbf{Patients} & \textbf{Taux r\'eadm.} & \textbf{Capture} \\
\midrule
Low & 10\,787 (55\,\%) & 7{,}3\,\% & 35{,}2\,\% \\
Medium & 8\,441 (43\,\%) & 15{,}4\,\% & 57{,}8\,\% \\
High & 545 (2{,}8\,\%) & 28{,}8\,\% & 7{,}0\,\% \\
\bottomrule
\end{tabular}
\end{center}

Le groupe HIGH a un taux de r\'eadmission $3{,}9\times$ plus \'elev\'e que le groupe LOW.

\subsection{SHAP --- Facteurs modifiables (ODS-B)}

SHAP (SHapley Additive exPlanations) d\'ecompose la pr\'ediction de chaque patient en contributions par feature. On identifie les facteurs \textbf{modifiables} --- ceux sur lesquels le patient peut agir :

\begin{center}
\small
\begin{tabularx}{\textwidth}{@{}lrX@{}}
\toprule
\textbf{Feature} & \textbf{$\overline{|\text{SHAP}|}$} & \textbf{Action propos\'ee} \\
\midrule
\texttt{num\_medications} & 0{,}093 & Utilisez un pilulier pour organiser vos m\'edicaments \\
\texttt{insulin} & 0{,}063 & Prenez votre dose d'insuline comme prescrit \\
\texttt{emergency\_cap} & 0{,}044 & Appelez votre m\'edecin avant d'aller aux urgences \\
\texttt{a1c\_tested} & 0{,}035 & Planifiez un test HbA1c \\
\texttt{change} & 0{,}032 & Suivez les changements de m\'edicaments prescrits \\
\texttt{outpatient\_bkt} & 0{,}016 & Rendez-vous de suivi dans les 7 jours \\
\bottomrule
\end{tabularx}
\end{center}

Apr\`es correction, \texttt{a1c\_tested} est un vrai facteur (SHAP 0{,}035 au lieu de 0) : l'action ``Planifiez un test HbA1c'' est maintenant d\'eclench\'ee pour les patients concern\'es. Pour l'insuline, la valeur indiqu\'ee est celle de sa modalit\'e la plus influente.

Pour chaque patient \`a risque moyen ou \'elev\'e, on g\'en\`ere un plan personnalis\'e en ne gardant que les facteurs modifiables dont le SHAP $> 0$ (poussent vers la r\'eadmission).


% ─────────────────────────────────────────────────────────────────────────────
\section{Lab 4 --- Am\'elioration (Tuning)}

\subsection{Recherche d'hyperparamètres}

\textbf{RandomizedSearchCV} : 40 combinaisons al\'eatoires, \'evalu\'ees en 5 folds (\texttt{GroupKFold} sur \texttt{patient\_nbr}), soit 200 ajustements au total. La m\'etrique optimis\'ee est le \textbf{PR-AUC} (adapt\'ee aux classes d\'es\'equilibr\'ees).

Param\`etres explor\'es :

\begin{center}
\small
\begin{tabular}{@{}lll@{}}
\toprule
\textbf{Param\`etre} & \textbf{Plage} & \textbf{R\^ole} \\
\midrule
\texttt{n\_estimators} & 200 -- 800 & Nombre d'arbres \\
\texttt{learning\_rate} & 0{,}01 -- 0{,}15 & Vitesse d'apprentissage \\
\texttt{max\_depth} & 3 -- 7 & Profondeur max \\
\texttt{num\_leaves} & 15 -- 63 & Taille des arbres \\
\texttt{min\_child\_samples} & 20 -- 100 & R\'egularisation par feuille \\
\texttt{subsample} & 0{,}60 -- 0{,}95 & Fraction des donn\'ees \\
\texttt{colsample\_bytree} & 0{,}50 -- 0{,}95 & Fraction des features \\
\texttt{reg\_alpha} (L1) & 0 -- 1 & Sparsit\'e \\
\texttt{reg\_lambda} (L2) & 0{,}5 -- 3 & Lissage des poids \\
\bottomrule
\end{tabular}
\end{center}

\subsection{R\'esultats du tuning}

\begin{center}
\begin{tabular}{@{}lccc@{}}
\toprule
\textbf{Mod\`ele} & \textbf{PR-AUC} & \textbf{ROC-AUC} & \textbf{Brier} \\
\midrule
Lab 2 --- LightGBM (avant) & 0{,}202 & 0{,}653 & --- \\
Lab 4 --- Tun\'e (apr\`es) & \textbf{0{,}207} & \textbf{0{,}662} & --- \\
Lab 4 --- Tun\'e + calibr\'e & 0{,}194 & 0{,}659 & \textbf{0{,}097} \\
\bottomrule
\end{tabular}
\end{center}

Meilleur score en validation crois\'ee (GroupKFold) : PR-AUC = 0{,}199. Am\'elioration sur le test : $+2{,}5\,\%$ sur PR-AUC, $+1{,}5\,\%$ sur ROC-AUC. Le mod\`ele tun\'e est l\'eg\`erement meilleur tout en gardant une bonne calibration (Brier $\approx 0{,}097$).

\subsection{Niveaux de risque (mod\`ele tun\'e)}

\begin{center}
\begin{tabular}{@{}lrrr@{}}
\toprule
\textbf{Niveau} & \textbf{Patients} & \textbf{Taux r\'eadm.} & \textbf{Capture} \\
\midrule
Low & 11\,922 (60\,\%) & 7{,}4\,\% & 39{,}6\,\% \\
Medium & 7\,068 (36\,\%) & 16{,}1\,\% & 50{,}6\,\% \\
High & 783 (4\,\%) & 28{,}2\,\% & 9{,}8\,\% \\
\bottomrule
\end{tabular}
\end{center}

Le groupe \textbf{MEDIUM + HIGH} capture \textbf{60{,}4\,\%} des r\'eadmissions tout en ne repr\'esentant que 40\,\% des patients --- l'\'equipe soignante peut concentrer ses efforts sur 4 patients sur 10 et couvrir 6 r\'eadmissions sur 10.


% ─────────────────────────────────────────────────────────────────────────────
\section{Conclusion}

\begin{enumerate}
  \item \textbf{12 features suffisent} : le mod\`ele patient conserve 96{,}8\,\% du PR-AUC du mod\`ele complet (90{,}9\,\% apr\`es calibration).
  \item \textbf{La calibration est essentielle} : sans elle, les probabilit\'es ne sont pas fiables (Brier 0{,}22 $\to$ 0{,}097).
  \item \textbf{Le tuning am\'eliore mod\'er\'ement} : +2{,}5\,\% PR-AUC --- le choix des features compte plus que les hyperparam\`etres.
  \item \textbf{V\'erifier les donn\'ees avant le mod\`ele} : un seul bug de lecture (\texttt{None} lu comme \texttt{NaN}) rendait une des 12 features constante.
  \item \textbf{SHAP rend le mod\`ele actionnable} : chaque patient \`a risque re\c{c}oit des actions concr\`etes sur ses facteurs modifiables.
\end{enumerate}

\end{document}
